% Audit contribution. Uses the theorem/proof environments and bibliography of
% the parent paper. No preamble and no private macros are required.
\subsection{The ordinary mapping-class tower and the completion audit}
\label{sec:tower-audit}

\subsection{The exact compatible-pair definition}

The supplied EXIT\_TOWER equation (E.200) is the equation (5.1) in
Lochak's Definition~5.4 \cite{Lochak2025Aut}, with
$\tau _1=u$, $\tau _2=w$ and $k=c$.
Thus the definition of a compatible pair in that passage has been
transcribed correctly.  Proposition--Definition~5.6 proves that projection
of the compatible-pair group onto its first GT parameter is injective.
Its image consists exactly of the first parameters whose forced second
parameter satisfies the GT relations.  The assertion that this projection
is surjective is the additional conclusion claimed in EXIT\_TOWER E.28;
it is not a conclusion of Proposition--Definition~5.6.  Definition~5.8
imposes equality of the two parameters, and Proposition~5.9 identifies the
resulting relation (IV) with the framed six-move relation (R$'$).
These are literal comparisons of the equations in the supplied PDF,
printed pp.~50--52 and 58--59.

Theorem~5.10 in that PDF concerns the procongruence mapping-class group.
EXIT\_TOWER E.42, in contrast, explicitly uses full profinite completions
and supplies its own proposed comparison.  The actual map between these
completions is
\[
 \widehat\Gamma_{g,n}\longrightarrow\check\Gamma_{g,n}
 =\operatorname{im}\bigl(\widehat\Gamma_{g,n}
             \longrightarrow\operatorname{Out}(\widehat\pi_{g,n})\bigr).
\]
It is the quotient by the kernel of the indicated outer action.  Thus
transporting a theorem from the target of this map to its source requires
control of this kernel or a separate proof on the source.  The E.42 proof
attempts the latter, through its full transporter-completion lemmas.

\subsection{The coefficient and the local relations}

\begin{result}{A.6.1}{Proved}{The ordinary full profinite tower used in E.49}
The same-parameter relation proved in A.1 extends the standard
genus-zero action of every $A\in\widehat{\mathrm{GT}}$ to genuine
quilt-dependent automorphisms of the full profinite mapping-class
groups of all stable labelled surfaces.  Changes of quilt give the
explicit inner comparisons below.  The induced outer actions and
their groupoid form respect labelled sewing, self-sewing, ordinary
marked capping and forgetting marked points.  In particular this
supplies the actions and forgetful comparisons on
$\widehat\Gamma_{1,9}$ and $\widehat\Gamma_{1,12}$ used in E.49.
The construction uses the explicit Schreier coefficient $c_2$;
no identification with the symmetrizer extension of the Kummer
cocycle is assumed.  It is independent of the whole-arrow
classification E.42.
\end{result}

Write $A=(\lambda,f)\in\widehat{\mathrm{GT}}$,
$\nu=(\lambda-1)/2$, and let $c=c_2(A)$ be the Schreier coefficient
of EXIT\_TOWER E.22.  The index-two subgroup of
$\widehat F(x,y)$ defined by $x\mapsto1$, $y\mapsto0$ in
$\mathbb Z/2$ has free profinite basis
\[
              x^2,\quad y,\quad xyx^{-1}.
\]
The abelianization of $f$ in this basis is $(0,c,-c)$.
For the GT automorphism
$\phi_A(x)=x^\lambda$,
$\phi_A(y)=f^{-1}y^\lambda f$, the subgroup is preserved and
its abelianization is multiplied by $\lambda$.  Indeed, $f$ belongs
to this subgroup, so the image of $y$ has class $\lambda[y]$.
Conjugation by $x^\lambda$ acts on its abelianization as conjugation
by $x$, since $\lambda$ is odd and $x^2$ belongs to the subgroup.
The three stated basis vectors consequently have scalar images.
The product convention $\phi_{AA'}=\phi_A\phi_{A'}$ has second
word $f_A\phi_A(f_{A'})$.  Taking its second Schreier coordinate
proves the actual identity
\[
 c_2(AA')=c_2(A)+\lambda(A)c_2(A').                 \tag{TA.1}
\]
This proof uses no symmetrizer relation.

This point matters in comparing with Nakamura--Schneps
\cite[Definition~5.1 and Proposition~5.3]{ET_NS2000}.
Their Definition~5.1 introduces a coefficient through the symmetrizer;
Proposition~5.3 identifies it with $-A_F(-1,1)$ on their subgroup
satisfying (III$'$).  The end of their \S5 also describes
$-A_F(-1,1)$ as an extension over all GT.  The extension by itself
does not identify those two constructions on every GT element.
The calculation here uses the explicitly defined Schreier coefficient
and (TA.1), so it needs no such identification.

The same-parameter relation proved in the preceding part of this audit is
\[
 f(u,w^4)=w^{8c}f(u^2,w^2)u^{4c}(uw)^{-6c}
                         \quad\hbox{in }\widehat B_3.       \tag{TA.2}
\]
We now record exactly how this relation supplies the remaining local
relation of the ordinary tower.  In $B_4$ put
\[
\begin{gathered}
 a=\sigma_1,\quad b=\sigma_2,\quad d=\sigma_3,\qquad
 X=a^2,\quad Y=b^2,\quad Z=d^2,\quad
 U=dba^2b^{-1}d^{-1},\\
 V=(aba)^2,\quad W=(dbd)^2,\quad\omega=(abd)^4,\qquad
 e_1=V^2,\quad e_2=dbab^{-1}d^{-1},\quad e_3=W^2.
\end{gathered}
\]
The elementary identities used below are
\[
\begin{array}{c|cccc}
(p,q)&(a,W)&(e_2,W)&(e_2,V)&(d,V)\\ \hline
(pq)^2&\omega Z&\omega Y&\omega Y&\omega X
\end{array},\qquad UWV=\omega Y,\qquad e_2^2=U.       \tag{TA.3}
\]
Each row also satisfies $(pq)^2=(qp)^2$.  The map
$u\mapsto p$, $w^2\mapsto q$ therefore defines a homomorphism
from the mixed subgroup $\langle u,w^2\rangle$ of $B_3$.
Applying (TA.2) and antisymmetry gives
\[
\begin{aligned}
 f(e_3,a)&=\omega^{2c}Z^{2c}X^{-2c}f(W,X)W^{-4c},\\
 f(e_2,e_3)&=\omega^{-2c}W^{4c}f(U,W)U^{2c}Y^{-2c},\\
 f(e_1,e_2)&=\omega^{2c}Y^{2c}U^{-2c}f(V,U)V^{-4c},\\
 f(d,e_1)&=\omega^{-2c}V^{4c}f(Z,V)Z^{2c}X^{-2c}.
\end{aligned}                                                   \tag{TA.4}
\]
Let $R$ be the word
\[
\begin{split}
R={}&f(e_3,a)d^{-8c}f(b^2,d^2)W^\nu
       f(e_2,e_3)e_2^{2\nu}f(e_1,e_2)b^{-2\nu}\\
 &\hspace{12mm}\cdot f(a^2,b^2)a^{8c}V^\nu
                    f(d,e_1)\omega^{-\nu}.
\end{split}                                                     \tag{TA.5}
\]
Substitute (TA.4).  The four powers of the central element $\omega$
cancel.  The commutations needed to put the remaining factors in order
are precisely
\[
 [X,Z]=1,\quad [Y,U]=[Y,V]=[Y,W]=1,\quad
 [V,X]=1,\quad[W,Z]=1.
\]
They give
\[
\begin{split}
R={}&(XZ)^{-2c}f(W,X)f(Y,Z)
 [W^\nu f(U,W)U^\nu f(V,U)V^\nu]\\
 &\hspace{9mm}\cdot Y^{-\nu}f(X,Y)f(Z,V)(XZ)^{2c}\omega^{-\nu}.
\end{split}
\]
The central-product hexagon and $UWV=\omega Y$ give
\[
 W^\nu f(U,W)U^\nu f(V,U)V^\nu=f(V,W)Y^\nu\omega^\nu.
\]
Consequently $R=(XZ)^{-2c}P(XZ)^{2c}$, where
\[
 P=f(W,X)f(Y,Z)f(V,W)f(X,Y)f(Z,V).
\]
Use $V=Xx_{13}Y$ and $W=Yx_{24}Z$.  A central factor commuting
with both arguments can be deleted from an argument of a commutator
proword.  Thus
\[
P=f(Yx_{24},X)f(Y,Z)f(Xx_{13},x_{24}Z)f(X,Y)f(Z,x_{13}Y).
\]
The ordinary pentagon identifies its middle three factors with
$f(X,Yx_{24})f(x_{13}Y,Z)$, and antisymmetry cancels the first
and last factors.  This proves $R=1$ in the full group
$\widehat B_4$.  The three-chain homomorphism to the two-boundary
torus sends $\omega$ to the product of its two boundary twists;
the boundary correction in (TA.5) is therefore retained.

All finite braid identities in (TA.3), the displayed commutations,
and the formulas for $V,W$ have also been verified by the faithful
Artin representation $B_4\hookrightarrow\operatorname{Aut}(F_4)$
in the accompanying script
\texttt{check\_tower\_braid\_identities.py}.
This certificate checks exact reduced words, rather than their images
in a selected finite group.

\subsection{Why the three-S relation and the group law use only \texorpdfstring{$c_2$}{c2}}

The three-S calculation can be made independently of any choice of
extension of the Kummer coefficient.  On a one-holed torus, write
$a,b$ for twists about curves of intersection one,
$d=bab^{-1}$, and $\sigma=aba$.  The braid identities give
\[
 b^2a^2d^2=\sigma^2,
 \qquad(aba)^2=(bdb)^2=(dad)^2,
 \qquad D_\partial=\sigma^4.
\]
For the three adjusted S-edges, the only coefficient factors that
remain after adjacent cancellations are $a^{-8c}$ at the start
and $a^{8c}$ at the end.  Their product is
\[
\begin{split}
 a^{-8c}f(b^2,a^2)\sigma^{2\nu}
 b^{2\nu}f(d^2,b^2)\sigma^{2\nu}
 d^{2\nu}f(a^2,d^2)a^{8c}\sigma^{2\nu}.
\end{split}                                                     \tag{TA.6}
\]
The central-product hexagon with $(x,y,z)=(b^2,a^2,d^2)$ is
\[
 f(b^2,a^2)b^{2\nu}f(d^2,b^2)d^{2\nu}
 f(a^2,d^2)a^{2\nu}=\sigma^{2\nu}.
\]
Since $\sigma^2$ is central, (TA.6) equals
\[
              \sigma^{8\nu}a^{-2\nu}
                 =D_\partial^{2\nu}a^{-2\nu}.               \tag{TA.7}
\]
The final quilt differs by $D_\partial^{-1}a$ from the initial
one; its adjustment weight is $D_\partial^{-2\nu}a^{2\nu}$,
which cancels (TA.7).  The equality holds for every value of $c$.
It does not use (III$'$) or the symmetrizer.

For completeness, the product law on an S-edge is just as explicit.
Let the adjustment integer be $N$ and put
\[
 r=N\nu-8c,\quad g=f(b^2,a^2),\quad h=a^r g,
 \quad t_A=h b^{8c}\sigma^{2\nu}.
\]
The action in the source quilt satisfies
\[
 \Phi_A(a)=a^\lambda,\quad
 \Phi_A(b)=h b^\lambda h^{-1},\quad
 \Phi_A(\sigma^2)=\sigma^{2\lambda}.
\]
The last equality is the standard central-full-twist formula in
$\widehat B_3$, also obtained by applying these first two formulas
to its central generator.  For primed parameters, direct substitution
and cancellation of $h^{-1}h$ give
\[
\begin{split}
 \Phi_A(t_{A'})t_A
 ={}&a^{r+\lambda r'}g\,
 f'\bigl(b^{2\lambda},g^{-1}a^{2\lambda}g\bigr)\\
 &\hspace{12mm}\cdot b^{8(c+\lambda c')}
                   \sigma^{2(\nu+\lambda\nu')}.
\end{split}                                                     \tag{TA.8}
\]
The two prowords are exactly the GT product evaluated at
$(b^2,a^2)$.  Equation (TA.1) and
$\nu(AA')=\nu(A)+\lambda(A)\nu(A')$ identify the exponents.
Thus (TA.8) is the required cocycle identity on this edge.
No assertion about the equality of two Kummer extensions has entered.

The other local cells have their stated elementary checks.  A-edge
backtracking uses $f(x,y)f(y,x)=1$.  S-edge backtracking cancels
$b^{8c}b^{-8c}$ and the inverse prowords, leaving
$a^{N\nu}\sigma^{4\nu}=a^{N\nu}D_\partial^\nu$,
which is its recorded quilt adjustment.  A three-A cell is the
central-product hexagon with the lantern relation.  A five-A cell is
the pentagon, with zero boundary exponents; the boundary twists
inject into the abelianization of the bordered five-point group.
Disjointly supported cells commute.  Equation (TA.5) is the six-AS
cell, including both boundary twists.  These are all the cells in
the simply connected pants complex of Hatcher's theorem
\cite[Theorem~8.2]{ET_NS2000}.
The NS calculations used for the geometric signs are
\cite[Definition~7.2, Lemma~7.3 and Claim~8.3]{ET_NS2000};
the algebraic verifications above identify
their complete coefficient dependence.  In particular no further
pentagon loop is hidden in the three-S or product-law step.

\subsection{The actual group action and ordinary forgetting}

Use quilted pants decompositions on a surface with nonempty
parametrized boundary, with the two external seam endpoints fixed.
Over a fixed pants decomposition, the half-adjustments form the free
integer lattice with one coordinate for each internal and external
circle.  To see this directly, cut every pant into two hexagons.
At a collar there is an endpoint interchange and an integral winding,
recorded by one integer in half-turn units.  When two cuffs are glued,
the common endpoint motion identifies one primitive diagonal in their
two collar coordinates.  With $v$ pants, $e$ internal circles and
$n$ external circles, the resulting lattice has rank
$3v-e=e+n$.  Its coordinates are the stated circle coordinates.
Freeness follows by squaring a putative stabilizing half-adjustment:
the resulting multitwist fixes the seam cellulation relative to the
external boundary.  Its edges can be fixed successively and each
hexagon then has the Alexander isotopy to the identity.  Independence
of the multitwists makes every coordinate zero.

Give an adjustment $D_a^{N/2}$ weight $D_a^{N\nu}$, an A-edge
replacing $a$ by $b$ weight $a^{N\nu}f(b,a)$, and an S-edge weight
\[
        a^{N\nu-8c}f(b^2,a^2)b^{8c}(aba)^{2\nu}.
\]
The preceding cell calculations show that a lifted pants cell has
weight inverse to its terminal quilt adjustment.  Any pants loop is
a product of these cells and backtracking.  The free adjustment
lattice then makes a closed quilt path have weight one.  Thus path
weights define $T_A(Q,Q')$ with
\[
 T_A(Q,Q'')=T_A(Q,Q')T_A(Q',Q''),\qquad
 T_A(hQ,hQ')=hT_A(Q,Q')h^{-1}.                              \tag{TA.9}
\]
For a boundary-preserving mapping class or morphism
$\gamma:S_G\to S_H$, set
\[
       \mathcal F_A(\gamma)
       =\gamma T_A(Q_G,\gamma^{-1}Q_H)^{-1}.                \tag{TA.10}
\]
Inserting (TA.9) in (TA.10) proves preservation of composition.
This is a homomorphism into a profinite group on each vertex group;
its composition with any finite quotient has finite-index kernel,
so it extends to the full profinite completion.  The edge cocycle,
proved on half-adjustments, A-edges, and explicitly in (TA.8) on
S-edges, extends along paths and proves
$\mathcal F_A\mathcal F_{A'}=\mathcal F_{AA'}$.
In particular $\mathcal F_{A^{-1}}$ is its inverse.

A path of quilts on a subsurface remains a path of the same local type
after sewing.  The two glued boundary twists both map to the positive
seam twist and their exponents add.  Each weight therefore maps to the
weight of the sewn path.  This proves sewing and self-sewing compatibility
on actual morphisms.  Relabelling follows from the second equality
in (TA.9).  The cylinder is assigned the automorphism
$\tau\mapsto\tau^\lambda$, in agreement with this same boundary rule.

We give the specific forgetting comparison needed for the main
finite-subgroup argument.  Fill an external boundary $\delta$ by a
disc.  The adjacent pant, with other cuffs $c,d$, becomes an annulus.
Join the two seam endpoints on $\delta$ by the unique disc arc relative
to those endpoints, and remove that annulus.  The two remaining seam
arcs are stitched into the adjacent pants.  On collar coordinates this
operation kills the $\delta$ coordinate and adds the $c,d$ coordinates.
The rule follows by lifting each collar to
$\mathbb R\times[0,1]$: concatenating the strips adds their signed
displacements, including half-integer displacements.  In particular
\[
 D_\delta\mapsto1,\qquad D_c,D_d\mapsto D,
 \qquad D_c^{1/2}Q,D_d^{1/2}Q\mapsto D^{1/2}\kappa Q.
                                                               \tag{TA.11}
\]
An A-move supported away from the cap keeps its weight.  An A-move
whose four-holed support contains the cap becomes a recutting of the
same three-holed quilt: its two replacement twists are now boundary
twists, hence commute, and its commutator proword is one.  Its only
remaining weight is the projected adjustment in (TA.11).  This
recutting adds no winding, since the stitched arcs are the same arcs
before and after regrouping the hexagons.  An S-move keeps its support
unless the source was the entire one-holed torus.  That exceptional
last cap is the central quotient discussed below.  These cases exhaust
the possible local supports.  Therefore
\[
       \kappa_*T_A(Q,Q')=T_A(\kappa Q,\kappa Q')              \tag{TA.12}
\]
whenever a boundary remains.  Equations (TA.10) and (TA.12) give the
actual commuting capping diagram in these coordinates.  Multiple
caps stitch the same collar strips; their order gives the same map.

At the group level, unmarked capping is a quotient whose discrete
kernel is generated by the capped boundary twist and disc-pushing
products $D_cD_d^{-1}$.  Here $c,d,\delta$ bound the pant following
a simple based loop of the filled surface.  This kernel description
comes from the evaluation fibration for an oriented embedded disc:
the fibre rotation gives $D_\delta$, and a loop of its centre gives
the indicated bounding-pair product.  Simple based loops generate
the surface group.  Completion of a quotient is right exact, so the
completed kernel is the closed normal closure of these generators.
No injectivity of the completed disc-pushing group is required.
A single quilt path to a pants decomposition containing both $c,d$
has transport $t$, and the twist formula gives
\[
 \Phi_A(D_cD_d^{-1})=t(D_cD_d^{-1})^\lambda t^{-1},\qquad
 \Phi_A(D_\delta)=D_\delta^\lambda.                         \tag{TA.13}
\]
The same formula for $A^{-1}$ proves preservation with equality of
the closed kernel.  Hence the automorphism descends to the capping
quotient.  Marked capping has only the boundary-twist kernel.
Blow every puncture up to a labelled boundary circle and cap by marked
discs.  Forgetting a puncture is unmarked capping at its blown-up circle;
(TA.12) identifies the descended action with the already constructed
action on the target.  This proves the usual forgetting compatibility
on $\widehat\Gamma_{g,n}$ with nonempty marking set, including all
$\widehat\Gamma_{1,n}$ needed for the finite detector.

For the last marking in genus at least two use the same quotient from
one boundary, with the disc-pushing generators (TA.13).  The identity
component of the closed-surface diffeomorphism group is contractible,
so the same evaluation sequence supplies the kernel.  For a torus,
the map is the explicit quotient
\[
          \widehat B_3\longrightarrow
          \widehat{\operatorname{SL}_2(\mathbb Z)},
             \qquad \ker=\overline{\langle(ab)^6\rangle}.
\]
Its kernel is the boundary twist and is preserved.  Marked and unmarked
last capping give this same group quotient.  Thus the construction
includes closed targets.  A change of starting quilt changes each
representative by the inner automorphism given in (TA.9); the resulting
outer actions and forgetting comparisons are independent of that choice.

This is the ordinary full-profinite group-action bridge needed in the
main argument.  It is proved by the local formulas and explicit quotient
maps; it does not require the additional whole-exit-category rigidity
assertion E.42.

\subsection{The E.9 and E.28 geometric model}

The orbifold disk with four order-two cone points has group
$D_4=C_2(x_1)*C_2(x_2)*C_2(x_3)*C_2(x_4)$.
Its parity kernel $K_4$, with all four generators odd, is free of
rank three.  It is the fundamental group of the double cover branched
at the four cone points, a genus-one surface with two boundary circles.
The two cusp loops after marking these circles are represented by
$c=x_1x_2x_3x_4$ and its conjugate by $x_1$.
The Artin braid action on this orbifold group has
$\omega_4=(\sigma_1\sigma_2\sigma_3)^4$ acting by
$\operatorname{Inn}(c^{-1})$ in the declared right-action convention.
Thus the diagonal central element $(c,\omega_4)$ acts trivially in
the universal-curve quotient.  Its quotient gives the Birman extension
\[
 K_4\longrightarrow
 (K_4\rtimes B_4)/\langle(c,\omega_4)\rangle
       \longrightarrow B_4/\langle\omega_4\rangle.
\]
The base is the two-pointed torus mapping-class group in this
branched-cover realization; the middle term adds the moving third
point.  This identifies the morphisms in E.9, not just their ranks.

Cutting the lift of the arc joining $x_1,x_2$ leaves the orbifold group
\[
 B=\langle p=x_1x_2\rangle*C_2(x_3)*C_2(x_4).
\]
Its parity kernel $L$ has free basis
$p$, $q_0=x_3px_3$, $r=x_3x_4$.  Its two old cusps are
$c=pr$ and $c'=q_0r^{-1}$.  The branch-preserving cut stabilizer,
modulo the twist of the cut curve, is the pure five-point sphere group;
forgetting the moving third torus point becomes the usual
$\Gamma_{0,5}\to\Gamma_{0,4}$ map with kernel $L$.
For the topology assertion, $B$ is a free factor in the index-two
subgroup of $D_4$ making $x_1,x_2$ odd and $x_3,x_4$ even;
$L$ has index two in $B$.  Retraction to that free factor and the
finite-index comparison give the full induced profinite topology.

The cabling identities invoked by E.28 are actual strand-doubling
identities: doubling the first strand sends
\[
 x_{14}\longmapsto x_{15}x_{25},\qquad
 x_{24}\longmapsto x_{35},\qquad
 x_{34}\longmapsto x_{45}.
\]
They and the orbifold full-twist identity have been checked by exact
word reduction in the accompanying braid script.  Standard GT strand
cabling supplies the compatible action on these maps.  No failure was
found in these particular E.9/E.77 geometric identifications.  The main
same-parameter reconstruction can nevertheless proceed directly from
(TA.2), without using the global involution construction E.28.

\subsection{Full completion: the supported diagonal calculation and the missing files}

Here is the finite-quotient calculation behind the diagonal part of E.29.
Put $C=\operatorname{Stab}_\Gamma(\sigma)$,
$\Delta=\langle\tau_c:c\in\sigma\rangle$ and $H=C/\Delta$.
The multitwist subgroup has its full induced topology in $\Gamma$
(the first injection of \cite[Theorem~4.5]{ET_Boggi2014}
is the cited input).
Given $N\triangleleft C$ of finite index, choose
$U\triangleleft\Gamma$ of finite index with $\Delta\cap U\subseteq N$.
Choose $V\triangleleft H$ of finite index contained in the image of
$N\cap U$.  For $c\in C\cap U$ with $p(c)\in V$, choose
$n\in N\cap U$ with $p(n)=p(c)$.  Then
$cn^{-1}\in\Delta\cap U\subseteq N$.  Consequently
\[
       C\cap U\cap p^{-1}(V)\subseteq N.
\]
This proves that the diagonal $C\to\Gamma\times H$ has its full
profinite topology, and hence gives the embedding
\[
          \widehat C\hookrightarrow\widehat\Gamma\times\widehat H.
                                                               \tag{TA.14}
\]
The completed crossing is $\widehat\Gamma/\widehat\Delta$.
For an element in its two-endpoint stabilizer, lift its
$\widehat H$ coordinate to $\widehat C$; the residual
$\widehat\Gamma$ coordinate belongs to $\widehat\Delta$.
It too comes from $\widehat C$.  Thus (TA.14) is the entire
crossing stabilizer.  This argument proves precisely that diagonal
claim, without appealing to a normalizer theorem for a general
multicurve in $\widehat\Gamma$.

There is, however, a literal missing-proof issue in the supplied archive.
The registry labels the following results ``Proved'' and names these
proof files, but none of the five files is present:
\[
\begin{array}{c|l}
 E.29&\texttt{sections/E\_arrow\_topology.tex}\\
 E.32&\texttt{sections/E\_local\_generizations.tex}\\
 E.33&\texttt{sections/E\_rational\_tail\_normalizers.tex}\\
 E.35&\texttt{sections/E\_sewing\_extraction.tex}\\
 E.38&\texttt{sections/E\_adjacent\_presentation.tex}.
\end{array}
\]
The first unresolved supplied-proof invocation in the proof of E.39 is
its opening use of E.32 to identify the completed adjacent incidence
biset and its full stabilizer.  Later it invokes E.38 for all relations.
E.40 and E.42 additionally invoke E.33 and the full flag and equalizer
parts of E.32, and the sewing extraction invokes E.35.
Accordingly the supplied archive alone does not contain a complete
written proof of its whole-arrow classification E.42.
Section~\ref{sec:completion-repair} supplies the needed full
finite-quotient, flag, normalizer, square-equalizer and sewing
arguments, including graph-automorphism coherence, and proves A.6.
The diagonal argument here is its first local calculation.
This records a repaired dependency gap in the original archive;
the ordinary forgetting bridge A.6.1 is already independent of
these missing files.

\subsection{A precise correction to the claimed Galois normalization}

Lochak's Remark~5.2(2), printed p.~59 \cite{Lochak2025Aut},
asserts disjointness of the
extensions generated by roots of unity and by all roots of $2$, and
uses it to allow arbitrary cyclotomic and Kummer parameters.
There is a direct algebraic obstruction.  For $\zeta=\zeta_8$,
\[
         (\zeta+\zeta^{-1})^2=2,
 \qquad\sqrt2\in\mathbb Q(\mu_\infty)
                      \cap\mathbb Q(2^{1/n}:n\geq1).
\]
With the positive-radical convention
$\sigma(2^{1/n})/2^{1/n}=\zeta_n^{\rho_2(\sigma)}$, an element
$\sigma\in G_{\mathbb Q}$ therefore satisfies
\[
  \rho_2(\sigma)\bmod2
       =\frac{\chi(\sigma)^2-1}{8}\bmod2.                    \tag{TA.15}
\]
Indeed, the four possible residues $\chi=1,3,5,7\pmod8$
send $\zeta+\zeta^{-1}$ respectively to
$\sqrt2,-\sqrt2,-\sqrt2,\sqrt2$.
The forced Kummer parities are $0,1,1,0$.
In particular $(\chi,\rho_2)=(1,1)$ is impossible.
The finite calculation in $\mathbb Q[\zeta]/(\zeta^4+1)$ is included
in the braid-and-coefficient script.  This corrects the arbitrary
normalization assertion.  It does not rule out a normalization for
parameter pairs satisfying their actual compatibility constraints.

\subsection{No parameter is lost on the rational-tail subgroup}

For the parameter comparison in E.37 put
$G=\widehat B_4/\langle\omega_4\rangle$,
$Q=\widehat{\operatorname{SL}_2(\mathbb Z)}$, and
$J=\langle a,b\rangle\simeq\widehat B_3$.  The explicit completed
Birman sequence E.63 is
$1\to K=\widehat F(X,Y)\to G\to Q\to1$.
Writing $z=(ab)^3$, its actual cusp identification is
$z^2=[X,Y]^{-1}$, so $J\cap K=I=\langle[X,Y]\rangle$ and
$J\to Q$ is surjective.  In the index-two free subgroup with basis
$A=X$, $B=YXY^{-1}$, $C=Y^2$, the word $[X,Y]=AB^{-1}$ is
primitive.  Thus $I$ has its full topology.  Right exact completion
of the two central sequences gives the injection
$\widehat B_3\hookrightarrow G$, with that exact intersection.
A normalizer of $J$ may be multiplied by an element of $J$ to enter
$K$; it then normalizes $I$.  The peripheral finite-cover argument
(TC.4) puts it in $I$.  Hence $N_G(J)=J$.

We can also prove directly the required injectivity
$\widehat{\mathrm{GT}}\to\operatorname{Out}(\widehat B_3)$.
Suppose the standard action of $(\lambda,f)$ is inner.  Its action
on $B_3^{\rm ab}=\widehat{\mathbb Z}$ forces $\lambda=1$.
An inner conjugator fixes $a$, hence belongs to
$C_{\widehat B_3}(a)=\langle a,z\rangle$, as proved by the
primitive-generator argument in E.26.  Modulo $z$ write it $a^t$.
The pure quotient is the free profinite group
$F=\langle x=a^2,y=b^2\rangle$ of index six.  On its
abelianization, an odd power of $a$ sends
$[x]\mapsto[x]$, $[y]\mapsto-[x]-[y]$; an even power acts
trivially.  The standard multiplier-one GT action acts trivially,
so $t=2s$.  Its equality on $y$ now says
$f x^s\in C_F(y)=\langle y\rangle$.
Since $f$ has zero abelianization, comparison in
$F^{\rm ab}=\widehat{\mathbb Z}^2$ forces $s=0$ and the remaining
y-power to be zero.  Thus $f=1$.

Minamide--Nakamura's classification\cite[Theorem~4.3(i), Corollary~C]{ET_MN2022} writes the outer class of any
automorphism of $G$ as the class of a unique genuine standard GT braid
action.  The adapted-quilt action in E.37 and that braid representative
both preserve $J$.  Their ambient inner conjugator consequently lies
in $N_G(J)=J$; restriction gives an inner conjugacy in $J$ itself.
The just-proved injectivity identifies their GT parameters exactly.
The restriction has no additional kernel.  Finally, restriction to
$\widehat P_4/\langle\omega_4\rangle$ retains that parameter;
an ambient inner representative can insert only the induced branch
permutation in $\mathfrak S_4$.  The given branch ordering fixes this
finite ambiguity.  This is the pinned parameter comparison used by E.37.
